% ===========================================================================
%  E08 — Ferramentas do professor
% ===========================================================================
\chapter{E08 \textperiodcentered{} Ferramentas do professor}
\label{ep:e08}

\begin{fichaepico}
\begin{campos}
  \item[Objetivo] Tirar do professor o trabalho repetitivo do planejamento e
    cuidar de quem cuida da turma.
  \item[Para quem] Professor.
  \item[Valor] São três produtos que a rede pode contratar separadamente do
    Entre Redes: Gerador de Plano de Aula, Planejamento de Ensino e Como você
    está?.
  \item[Estado] \pronto.
\end{campos}
\end{fichaepico}

\begin{tabularx}{\linewidth}{@{}L{16mm}X L{20mm}@{}}
\toprule
\headrow \thd{ID} & \thd{Feature} & \thd{Estado}\\
\midrule
\id{F08.1} & Gerador de Plano de Aula & \pronto\\
\zebra \id{F08.2} & Planejamento de Ensino & \pronto\\
\id{F08.3} & Como você está? & \pronto\\
\bottomrule
\end{tabularx}

\begin{nota}
\textbf{Dois planejadores, dois públicos.} O Gerador de Plano de Aula é para a
rede que já tem currículo próprio e material didático adotado: ele parte desse
material. O Planejamento de Ensino é para quem parte da BNCC e quer descer até
a aula de cada semana. Os dois convivem; um não substitui o outro.
\end{nota}

% ---------------------------------------------------------------------------
\feature{F08.1}{Gerador de Plano de Aula}{\pronto}

\begin{historia}{H08.1.1}{Gerar o plano anual a partir do material da rede}{\pronto}
\hu{Professor}{gerar um plano anual da minha disciplina e série}{não começar
o ano com a página em branco}
\aceite
\begin{itemize}
  \item \enquote{Meus planos de aula} \seta{} \enquote{Novo plano de aula}:
    escolho disciplina e série.
  \item Aparecem só as séries com material cadastrado pela rede; disciplinas
    sem material aparecem como \enquote{em breve}.
  \item A geração leva alguns minutos, com dicas na tela enquanto espero;
    posso fechar a tela, e ela continua.
  \item O plano começa pela fundamentação e pelos objetivos do ano, mostra o
    percurso do ano numa tabela (a pergunta, a aprendizagem central e o produto de
    cada bimestre) e, em cada bimestre, traz a pergunta norteadora, os conteúdos, as
    habilidades, as estratégias, as formas de avaliação e o produto de síntese.
  \item As formas de avaliação dizem o que o estudante mostra que aprendeu, e
    cada uma se liga a uma habilidade do bimestre.
  \item Um erro aparece como uma frase que eu entendo.
  \item Usa o meu limite diário de IA.
\end{itemize}
\pendencia{depende de a rede ter enviado o material de referência.}
\end{historia}

\begin{historia}{H08.1.2}{Editar e levar o plano para o Word}{\pronto}
\hu{Professor}{ajustar o plano gerado e exportar}{entregar à coordenação no
formato que ela pede}
\aceite
\begin{itemize}
  \item O plano abre como documento editável: salvar (Ctrl/Cmd+S), renomear,
    excluir.
  \item \enquote{Restaurar original} volta à versão que a IA gerou, que fica
    guardada intacta.
  \item Exportar em .doc (abre no Word).
  \item O plano é do professor: não fica público.
\end{itemize}
\end{historia}

% ---------------------------------------------------------------------------
\feature{F08.2}{Planejamento de Ensino}{\pronto}

\begin{historia}{H08.2.1}{Declarar a minha regência}{\pronto}
\hu{Professor}{informar o que eu leciono e a minha grade}{que o planejamento
tenha as datas reais do meu ano}
\aceite
\begin{itemize}
  \item Informo componente, série, ano letivo e data de início, grade de
    horários, minutos por aula e divisão do ano (bimestre, trimestre ou
    semestre).
  \item A carga horária sugerida vem da mediana de redes reais e aparece como
    sugestão, que posso mudar.
  \item Hora-aula não é hora-relógio: a conta respeita os minutos por aula.
\end{itemize}
\end{historia}

\begin{historia}{H08.2.2}{Distribuir as habilidades da BNCC no ano}{\pronto}
\hu{Professor}{colocar cada habilidade da BNCC num período do ano}{garantir
que nenhuma fique de fora}
\aceite
\begin{itemize}
  \item Cada habilidade vai para um período.
  \item No Ensino Médio, a tela avisa quando as habilidades são por área.
  \item Nenhuma habilidade fora do período entra no plano daquele período.
\end{itemize}
\end{historia}

\begin{historia}{H08.2.3}{Descer do plano anual à aula da semana}{\pronto}
\hu{Professor}{ter o plano anual e o plano de cada semana}{chegar à sala
sabendo o que fazer em cada aula}
\aceite
\begin{itemize}
  \item O plano anual é gerado e eu posso reescrevê-lo.
  \item As 40 semanas são planejadas uma por vez, com tema, objetivo, quatro
    momentos, recursos, avaliação, tarefa e adaptações.
  \item As datas vêm do calendário real: feriado remove a aula. A IA preenche
    o conteúdo, mas nunca decide quantas aulas cabem nem inventa código de
    habilidade.
  \item Marco \enquote{AULA CUMPRIDA} e escrevo observações.
\end{itemize}
\end{historia}

\begin{historia}{H08.2.4}{Exportar a semana, a quinzena ou o bimestre}{\pronto}
\hu{Professor}{baixar o planejamento no recorte que a coordenação
pede}{entregar sem reformatar}
\aceite
\begin{itemize}
  \item \enquote{Baixar esta semana}, \enquote{Esta e a próxima (quinzena)} e
    \enquote{Baixar o Nº Bimestre inteiro}, em .doc; ou imprimir.
  \item Semanas cuja geração falhou não entram no documento.
  \item Posso excluir o planejamento daquele componente e série.
\end{itemize}
\end{historia}

% ---------------------------------------------------------------------------
\feature{F08.3}{Como você está?}{\pronto}

Um acompanhamento curto do bem-estar do próprio professor: três frases por dia
de um questionário reconhecido (DASS-21, que olha depressão, ansiedade e
estresse), e sete dias fecham o retrato da semana. Leva menos de um minuto.

\begin{historia}{H08.3.1}{Responder hoje em menos de um minuto}{\pronto}
\hu{Professor}{responder três frases sobre como estou}{acompanhar o meu
bem-estar sem tomar tempo do dia}
\aceite
\begin{itemize}
  \item O cartão no Início mostra se já respondi hoje.
  \item \enquote{Responder hoje}: três frases numa escala e, se eu quiser, um
    texto livre; \enquote{Enviar respostas}; recebo uma mensagem de
    acolhimento.
  \item Uma resposta por dia.
  \item A tela diz com clareza que nenhuma pessoa da escola ou da rede lê o
    texto livre, e que uma IA lê.
  \item As frases do questionário nunca são reescritas.
\end{itemize}
\end{historia}

\begin{historia}{H08.3.2}{Ver o retrato da minha semana}{\pronto}
\hu{Professor}{ver o resultado ao completar sete dias}{perceber como estou ao
longo do tempo}
\aceite
\begin{itemize}
  \item \enquote{Falta N dias para o retrato da semana}; no sétimo dia,
    \enquote{Retrato da sua semana} com a pontuação de cada dimensão e a faixa
    correspondente.
  \item Um ciclo que fica 21 dias sem fechar expira, sem pontuação.
  \item \enquote{Seu histórico} mostra os ciclos concluídos e o que escrevi em
    cada dia.
\end{itemize}
\end{historia}

\begin{historia}{H08.3.3}{Receber apoio quando preciso}{\pronto}
\hu{Professor que está passando por um momento difícil}{receber orientação de
onde buscar ajuda}{não ficar sozinho com isso}
\aceite
\begin{itemize}
  \item Se as respostas ou o texto indicam risco, aparece uma mensagem
    \emph{fixa} com o CVV (188), o SAMU (192) e a orientação de procurar o
    CAPS ou a UBS.
  \item Nesse caso, a IA não é chamada.
\end{itemize}
\end{historia}

\begin{historia}{H08.3.4}{Só meu, e apagável}{\pronto}
\hu{Professor}{ter certeza de que ninguém da gestão vê o meu bem-estar, e poder
apagar tudo}{responder com sinceridade}
\aceite
\begin{itemize}
  \item Não existe tela de gestão, número agregado nem exportação desse dado.
  \item \enquote{Apagar todas as suas respostas} remove tudo, de todos os
    ciclos --- e funciona mesmo se a rede desligar o produto.
  \item O dado de bem-estar nunca entra no perfil portátil.
\end{itemize}
\end{historia}
